% 天文学（综述）
% license CCBYSA3
% type Wiki

本文根据 CC-BY-SA 协议转载翻译自维基百科\href{https://en.wikipedia.org/wiki/Astronomy?utm_source=chatgpt.com}{相关文章}。

\begin{figure}[ht]
\centering
\includegraphics[width=6cm]{./figures/ddcd5c7b339e766a.png}
\caption{欧洲南方天文台（ESO）的帕瑞纳尔天文台（Paranal Observatory）向银河系中心发射激光导星（Laser Guide Star）。} \label{fig_Astron_1}
\end{figure}
天文学是一门研究天体及宇宙中发生的各种现象的自然科学。它运用数学、物理学和化学来解释这些天体的起源及其整体演化。研究对象包括行星、卫星、恒星、星云、星系、流星体、小行星和彗星等；相关的现象则包括超新星爆发、伽马射线暴、类星体、耀变体、脉冲星以及宇宙微波背景辐射等。更广义地说，天文学研究的一切事物都起源于地球大气层之外。而宇宙学（Cosmology）则是天文学的一个分支，专门研究宇宙整体的起源、结构与演化。

天文学是最古老的自然科学之一。在有文字记载的早期文明中，人类就已开始系统地观测夜空。其中包括古埃及人、巴比伦人、希腊人、印度人、中国人、玛雅人，以及美洲的许多原住民。在古代，天文学的范畴十分广泛，涵盖了天体测量学、天体导航、观测天文学以及历法制定等多个领域。

现代专业天文学大体上分为观测天文学与理论天文学两个分支。  
观测天文学的核心任务是通过望远镜与其他仪器对天体进行观测，从而获取数据；然后利用物理学的基本原理对这些数据进行分析。  
理论天文学则侧重于建立计算机模型或解析模型，用以描述天体及其物理现象。  
这两个分支相辅相成：理论天文学旨在解释观测结果，而观测天文学则通过实验数据来验证理论预测。

天文学是少数几个业余爱好者仍然能够发挥积极作用的科学领域之一。这一点在瞬变天体事件（如新星爆发、彗星出现等）的发现与观测中尤为明显。许多重要的天文发现——包括新的彗星——都是由业余天文学家首先发现的。
\subsection{词源}
“天文学”（Astronomy）一词源自希腊语ἀστρονομία（astronomia），由 ἄστρον（astron，“星星”）与 -νομία（-nomia，来自νόμος nomos，“法则”或“规律”）组成，意为“研究天体的学科”。\(^\text{[1]}\)天文学不应与“占星学”（Astrology）混淆——后者是一种信仰体系，声称人类事务与天体位置之间存在关联。二者在起源上有共同渊源，但后来分道扬镳：天文学建立在物理学的科学基础之上，而占星学则缺乏物理学支持。\(^\text{[2]}\)
\subsubsection{“天文学”与“天体物理学”的用法}
在现代用法中，“Astronomy”（天文学）与“Astrophysics”（天体物理学）基本上是同义词。\(^\text{[3][4][5]}\)按照词典定义，天文学是“研究地球大气层之外的物体与物质及其物理和化学性质的科学”；\(^\text{[6]}\)而天体物理学是天文学的一个分支，专门研究“天体的行为、物理性质及其动力过程”。\(^\text{[7]}\)在某些情况下（例如 Frank Shu 所著的《The Physical Universe》（《物理宇宙》）导论中），“天文学”被用来指代该学科的定性研究，而“天体物理学”则强调以物理为导向的定量研究。\(^\text{[8]}\)某些领域（如天体测量学Astrometry）在这种意义上被视为“纯粹的天文学”，而非天体物理学。研究机构或大学的命名方式往往取决于该部门是否历史上隶属于物理学系。\(^\text{[4]}\)事实上，许多职业天文学家持有的是物理学学位，而非天文学学位。\(^\text{[5]}\)因此，在当代学术语境中，“天文学”与“天体物理学”常被互换使用。\(^\text{[3]}\)
\subsection{历史}
\subsubsection{史前时期}
\begin{figure}[ht]
\centering
\includegraphics[width=6cm]{./figures/2ec12ec85a389120.png}
\caption{尼布拉星象盘（约公元前 1800–1600 年），在一处可能具有天文功能的遗址附近被发现。  该星盘极有可能描绘了太阳或满月、弯月、昴星团（Pleiades）以及夏至与冬至，后两者以盘侧的金条表示。\(^\text{[9][10]}\) 星盘的上缘代表地平线与北方。\(^\text{[11]}\)} \label{fig_Astron_2}
\end{figure}
天文学最初的发展是由实际需求所推动的，例如用于农业活动的历法制定。在有文字记载之前，考古遗址如巨石阵（Stonehenge）就已显示出人类对天文观测的浓厚兴趣。\(^\text{[12]: 15 }\)另一类证据来自文物，例如尼布拉星象盘（Nebra Sky Disc）。这件青铜器被认为是一种天文历法，其设计将一年定义为十二个太阴月（共 354 天），并通过闰月来校正与太阳年的差异。星象盘上镶嵌有象征性的图案，被解释为太阳、月亮和星辰，其中包含一个由七颗星组成的星团。\(^\text{[9][13][14]}\)
\subsubsection{古典时期}
\begin{figure}[ht]
\centering
\includegraphics[width=6cm]{./figures/4dd8438ae26db3a1.png}
\caption{巴比伦平面星图（约公元前 7 世纪）。巴比伦天文学是早期的一种天文观测工具。其所使用的六十进制（sexagesimal）体系（例如 12、24、60、360 等）至今仍广泛用于计时与天体测量（astrometry）。\(^\text{[15]}\)} \label{fig_Astron_3}
\end{figure}
像埃及、两河流域、希腊、印度和中国等文明，在跨文化交流的影响下，纷纷建立了天文观测台，并发展出关于宇宙本质的思想，以及用于测时的历法与天文仪器。\(^\text{[16]}\)天文学早期的一个关键发展是巴比伦人开创了数学与科学天文学，为其他文明的天文传统奠定了基础。\(^\text{[17]}\)巴比伦人发现了月食以沙罗周期（saros cycle）为规律重复出现，该周期为 223 个朔望月。\(^\text{[18]}\)

继巴比伦人之后，古希腊与希腊化世界在天文学上取得了重大进展。希腊天文学致力于对天体现象寻找理性与物理的解释。\(^\text{[19]}\)在公元前 3 世纪，萨摩斯的阿里斯塔克斯（Aristarchus of Samos）估算了月球与太阳的大小与距离，并提出了一个以太阳为中心、地球与行星环绕其旋转的太阳中心模型（heliocentric model）。\(^\text{[20]}\)在公元前 2 世纪，喜帕恰斯（Hipparchus）计算了月球的大小与距离，并发明了已知最早的天文仪器之一，如星盘（astrolabe）。\(^\text{[21]}\) 他还观测到春分点与秋分点、夏至与冬至相对于恒星位置的微小漂移，这一现象即如今所知的岁差（precession）。\(^\text{[12]}\)喜帕恰斯还编制了包含1020颗恒星的星表，北半球大多数星座的命名均源自希腊天文学。\(^\text{[22]}\)约公元前 150–80 年的安提基特拉机械（Antikythera mechanism）是一种早期的模拟计算机，用于计算太阳、月亮与行星在特定日期的位置。在其之后，直到 14 世纪欧洲出现机械天文钟，才再次出现与之相当复杂的技术装置。\(^\text{[23]}\)  

在古希腊古典时期之后，天文学长期受地心宇宙模型（geocentric model）或托勒密体系（Ptolemaic system）主导，该体系以克劳狄乌斯·托勒密（Claudius Ptolemy）命名。  
他以 13 卷本的著作《天文学大成（Almagest）》（阿拉伯译名）总结了天文学知识，  
并在此后一千多年中成为主要的参考文献。\(^\text{[24]: 196 }\)在该体系中，地球被认为是宇宙的中心，太阳、月亮与恒星围绕地球旋转。\(^\text{[25]}\)虽然该理论最终被推翻，但在当时它提供了最精确的天体位置预测模型。\(^\text{[24]}\)
\subsubsection{后古典时期}
\begin{figure}[ht]
\centering
\includegraphics[width=6cm]{./figures/983767d1c3fc939d.png}
\caption{《天文学汇编》（Compilatio astronomica, 1493）中阿尔弗拉格努斯（Alfraganus）的画像。伊斯兰天文学家收集并翻译了印度、波斯与希腊的天文学著作，并在此基础上加入了他们自己的研究成果\(^\text{[26]}\)。} \label{fig_Astron_4}
\end{figure}
在中世纪的伊斯兰世界，天文学蓬勃发展。早在 9 世纪初，伊斯兰世界就建立了天文观测台\(^\text{[27][28][29]}\)。964 年，波斯穆斯林天文学家阿卜杜勒·拉赫曼·苏菲（Abd al-Rahman al-Sufi）在其著作《恒星之书》（Book of Fixed Stars）中首次描述了仙女座星系（Andromeda Galaxy）—— 这是本星系群（Local Group）中最大的星系\(^\text{[30]}\)。在 1006 年，埃及天文学家阿里·伊本·里德万（Ali ibn Ridwan）与中国天文学家共同观测到SN 1006 超新星，这是过去一千年中视星等最亮的恒星事件\(^\text{[31]}\)。伊朗学者比鲁尼（Al-Biruni）指出，与托勒密不同，太阳的远地点（apogee）并非固定，而是会移动的\(^\text{[32][33]}\)。阿拉伯天文学家还为许多恒星命名，这些阿拉伯星名至今仍在国际天文学命名中沿用\(^\text{[34]}\)。  

大津巴布韦（Great Zimbabwe）与廷巴克图（Timbuktu）的遗址\(^\text{[35]}\)可能曾设有天文观测设施\(^\text{[36]}\)。在后古典时期的西非，天文学家研究恒星的运行与季节的关系，并基于复杂的数学计算绘制星空图与行星轨道示意图\(^\text{[37]}\)。松盖帝国史学家马哈茂德·卡提（Mahmud Kati）在 1583 年记录了一次流星雨\(^\text{[38]}\)。  

在中世纪欧洲，理查德·沃林福德（Richard of Wallingford, 1292–1336）发明了首个天文钟（astronomical clock）——“矩形仪”（Rectangulus），可用于测量行星与其他天体之间的角度\(^\text{[39]}\)；他还设计了一个计算仪（equatorium）“阿尔比恩（Albion）”，可用于计算月球、太阳及行星的经度等天文数据\(^\text{[40]}\)。尼古拉·奥雷斯姆（Nicole Oresme, 1320–1382）讨论了地球自转的可能性证据\(^\text{[41]}\)；而让·布里丹（Jean Buridan, 1300–1361）提出了冲力理论（theory of impetus），用于解释天体等物体的运动\(^\text{[42][43]}\)。在长达六个多世纪的时间里——从中世纪晚期古典学术复兴直到启蒙运动时期——罗马天主教会为天文学研究提供了比其他任何机构都更多的经济与社会支持。教会的主要动机之一是确定复活节的日期\(^\text{[44]}\)。
\subsubsection{早期望远镜时期}
\begin{figure}[ht]
\centering
\includegraphics[width=6cm]{./figures/f6b1f99c08c2382d.png}
\caption{伽利略于 1610 年发表的划时代著作《星际信使（Sidereus Nuncius）》中，首次绘制的月球地形草图。} \label{fig_Astron_5}
\end{figure}
在文艺复兴时期，尼古拉·哥白尼（Nicolaus Copernicus）提出了日心说（heliocentric model）的太阳系模型\(^\text{[45]}\)。1610 年，伽利略·伽利莱（Galileo Galilei）观测到金星具有类似月球的相位变化，这一发现为日心模型提供了支持\(^\text{[12]}\)。大约在同一时期，约翰内斯·开普勒（Johannes Kepler）以定量的方式整理了日心体系\(^\text{[46]}\)。他分析了第谷·布拉赫（Tycho Brahe）长达二十年的精密观测数据，建立了一个系统，能够描述行星围绕太阳运动的细节\(^\text{[47]: 4 [48]}\)。虽然开普勒抛弃了哥白尼提出的匀速圆周运动，改以椭圆轨道来解释行星运动\(^\text{[12]}\)，但他并未能提出解释这些运动规律的理论基础\(^\text{[49]}\)。最终是艾萨克·牛顿（Isaac Newton）通过天体动力学与万有引力定律揭示了行星运动的原因\(^\text{[50]}\)。牛顿还发明了反射式望远镜\(^\text{[51]}\)，并与理查德·本特利（Richard Bentley）合作提出，恒星与太阳本质相同，只是距离更为遥远\(^\text{[47]}\)。  

新的望远镜也改变了人类对恒星的认识。1610 年，伽利略发现夜空中横贯天际的光带—— 银河（Milky Way）—— 其实是由无数颗恒星组成的\(^\text{[44]}\)。1668 年，詹姆斯·格雷戈里（James Gregory）比较了木星与天狼星的亮度，估计木星的距离超过 83,000 个天文单位\(^\text{[47]}\)。英国天文学家、首任“皇家天文学家”约翰·弗兰斯蒂德（John Flamsteed）编制了超过 3000 颗恒星的目录，但他的数据在 1712 年被违背意愿地出版\(^\text{[52]}\)。天文学家威廉·赫歇尔（William Herschel）制作了详细的星云与星团目录，并于 1781 年发现了天王星（Uranus），这是人类发现的第一颗新行星\(^\text{[53]}\)。弗里德里希·贝塞尔（Friedrich Bessel）在 1838 年发展了恒星视差（stellar parallax）技术，但因其应用极为困难，到 1900 年仅约 100 颗恒星的距离被成功测定\(^\text{[47]}\)。  

在 18 至 19 世纪间，欧拉（Leonhard Euler）、克莱罗（Alexis Claude Clairaut） 与达朗贝尔（Jean le Rond d'Alembert）对三体问题（three-body problem）的研究，  
使得人们能更准确地预测月球与行星的运动。这一研究又被拉格朗日（Joseph-Louis Lagrange）与拉普拉斯（Pierre Simon Laplace）进一步完善，使得可以通过扰动计算估计行星与卫星的质量\(^\text{[54]}\)。  

随着新技术的引入，天文学取得了重大进展，包括光谱仪（spectroscope）与天体摄影（astrophotography）的发明。1814–1815 年，约瑟夫·冯·夫琅和费（Joseph von Fraunhofer）在太阳与其他恒星的光谱中发现了约 574 条暗线\(^\text{[55][56]}\)。1859 年，古斯塔夫·基尔霍夫（Gustav Kirchhoff）将这些谱线解释为不同化学元素存在的结果\(^\text{[57]}\)。
\subsubsection{星系}
\begin{figure}[ht]
\centering
\includegraphics[width=6cm]{./figures/431d515c0910b0a2.png}
\caption{} \label{fig_Astron_6}
\end{figure}
在 18 世纪晚期，威廉·赫歇尔（William Herschel）绘制了从地球出发、向不同方向观测到的恒星分布图，并得出结论：宇宙由一片恒星盘组成，太阳位于其中心附近——这就是银河系（Milky Way）。在约翰·米歇尔（John Michell）证明恒星在本征光度上存在差异之后，以及赫歇尔自己利用更强大望远镜进行观测时发现恒星在所有方向上都有分布，天文学家们开始推测，那些模糊的螺旋星云（spiral nebulae）其实可能是遥远的“岛宇宙”（island Universes）\(^\text{[47]: 6 }\)。
\begin{figure}[ht]
\centering
\includegraphics[width=6cm]{./figures/d8de8f09567f6ed6.png}
\caption{艾萨克·罗伯茨（Isaac Roberts）于 1888 年拍摄的仙女座“大星云”（Great Andromeda "Nebula"）照片。\(^\text{[59][60]: 63}\) } \label{fig_Astron_7}
\end{figure}
直到 20 世纪，天文学家才最终证明，包括地球所在的银河系（Milky Way）在内的星系，实际上是由大量恒星组成的集合体。\(^\text{[61]}\)1912 年，亨丽埃塔·勒维特（Henrietta Leavitt）发现了造父变星（Cepheid variable stars），其光度变化具有明确的周期规律，可以用来确定恒星的真实光度，从而成为一种精确的宇宙距离测量工具。利用造父变星，哈洛·沙普利（Harlow Shapley）绘制出了第一张精确的银河系结构图。\(^\text{[47]: 7 }\)随后，埃德温·哈勃（Edwin Hubble）使用胡克望远镜（Hooker Telescope），在多个螺旋星云中识别出造父变星，  
并于 1922–1923 年间确凿地证明：仙女座星云（Andromeda Nebula）与三角座星云（Triangulum）等，实际上都是位于银河系之外的完整星系。这一发现首次证明：宇宙由无数个星系组成。\(^\text{[62]}\)
\subsubsection{宇宙学}
1917 年，阿尔伯特·爱因斯坦（Albert Einstein）发表了广义相对论，这标志着现代宇宙整体理论模型时代的开始。[63]1922 年，亚历山大·弗里德曼（Alexander Friedman）发表了简化的宇宙模型，展示了静态、膨胀和收缩三种可能的解。[47]: 13  1929 年，埃德温·哈勃（Edwin Hubble）发表了观测结果，表明所有星系都以与其距离成正比的速度远离地球，这一关系即现在所称的哈勃定律（Hubble’s law）。若宇宙正在膨胀，则这种关系是预期结果。[47]: 13 由此推论出宇宙曾经极其致密而炽热—— 即大爆炸（Big Bang）的思想，最早由乔治·勒梅特（Georges Lemaître）在 1927 年提出。[64]当时这一观点虽被讨论，但尚无实验证据支持。自 1940 年代起，人们开始研究高密度条件下的核反应速率，在 1940 年代末至 1950 年代初成功建立了大爆炸核合成模型（big bang nucleosynthesis）。随后在 1965 年，宇宙微波背景辐射（cosmic microwave background radiation）被发现，为大爆炸理论提供了确凿证据。[47]: 16   

理论天文学预测了黑洞（black holes）[65]与中子星（neutron stars）[66]的存在，  
它们被用于解释诸如类星体（quasars）[67]和脉冲星（pulsars）[68]等天体现象。  

空间望远镜（space telescopes）的出现使得科学家能够在大气层通常会阻挡或模糊的电磁波谱区间内进行测量。[69]而在 2015 年，LIGO 项目探测到了引力波（gravitational waves）的证据，[70][71]这为爱因斯坦广义相对论的预言提供了又一重实证支持。
\subsection{观测天文学}
\begin{figure}[ht]
\centering
\includegraphics[width=6cm]{./figures/29f2885a50af0910.png}
\caption{} \label{fig_Astron_8}
\end{figure}
观测天文学依赖于电磁辐射的多种波长，并根据观测所涉及的电磁波谱区域（electromagnetic spectrum）将天文学划分为不同的分支。[72] 以下将分别介绍这些子领域的具体信息。
\subsubsection{射电天文学}
\begin{figure}[ht]
\centering
\includegraphics[width=6cm]{./figures/62211f8ffc667fb1.png}
\caption{} \label{fig_Astron_9}
\end{figure}
射电天文学（Radio Astronomy）使用波长较长的辐射，主要在 $1\,\text{mm}$ 至 $15\,\text{m}$ 的范围内（对应频率从 $20\,\text{MHz}$ 到 $300\,\text{GHz}$），远远超出可见光谱范围。[73]通常不可见的氢气（Hydrogen）在 $21\,\text{cm}$（$1420\,\text{MHz}$）处产生一条可在射电波段观测到的谱线。[74]在射电波长下可被观测的天体包括：星际气体（interstellar gas）[74]、脉冲星（pulsars）[74]、快速射电暴（fast radio bursts）[74]、超新星（supernovae）[75]以及活动星系核（active galactic nuclei）[76]。
\subsubsection{红外天文学}
\begin{figure}[ht]
\centering
\includegraphics[width=6cm]{./figures/cb188a7a2b061a78.png}
\caption{} \label{fig_Astron_10}
\end{figure}
红外天文学探测波长比红色可见光更长、超出人类视觉范围的红外辐射。红外波段非常适合研究那些温度太低而无法发射可见光的天体，例如行星、恒星周盘（circumstellar disks）或被尘埃遮蔽的星云。红外的长波长可以穿透阻挡可见光的尘埃云，从而使天文学家能够观测到分子云中被包裹的年轻恒星，以及星系的核心。例如，“广域红外巡天探测器”（Wide-field Infrared Survey Explorer, WISE）的观测在揭示大量银河系原恒星（protostars）及其宿主星团方面尤为高效。[77][78]

除接近可见光波段的红外辐射外，大多数红外辐射在地球大气层中会被强烈吸收或被掩盖，因为大气本身会发出显著的红外辐射。因此，红外天文台必须建在地球上高而干燥的地区，或直接部署在太空中。[79]某些分子在红外波段具有强烈的辐射特性，这使得天文学家能够研究宇宙空间中的化学组成。[80]

詹姆斯·韦布太空望远镜（James Webb Space Telescope）利用红外辐射来探测极其遥远的星系。这些星系的可见光在数十亿年前发出，随着宇宙的膨胀而被红移至红外波段。通过研究这些遥远星系，天文学家希望了解最早星系的形成过程。[81]
\subsubsection{光学天文学} 
从历史上看，光学天文学（又称“可见光天文学”）是最古老的天文学形式。[82]早期的观测图像是手工绘制的。在19世纪末至20世纪大部分时间里，观测图像通过照相设备记录。现代的观测图像使用数字探测器，尤其是电荷耦合器件（CCD）进行采集并以现代介质记录。虽然可见光波段约为 $380$–$700\,\text{nm}$，[83]但同样的设备也可以用来观测近紫外和近红外辐射。[84]
\subsubsection{紫外天文学}
紫外天文学利用地球大气层无法透过的紫外波长，因此需要在高层大气或太空中进行观测。该领域最适合研究热辐射与来自炽热蓝色OB型恒星的光谱发射线，这些恒星在紫外波段极为明亮。[85]
\subsubsection{射线天文学}

